\section{GRPO with look-ahead}
\label{sec:method}

Throughout, a \emph{turn} is one utterance by either speaker, a \emph{prefix} $c$ is a transcript
ending on a patient turn, $\oplus$ appends turns to a transcript, and $\pi_n$ is the therapist
policy at the start of training iteration~$n$ ($\pi_1$ is the base model). The patient simulator
$P$ and the oracle $O$ are fixed LLMs; $O$ returns a scalar MI score for a transcript.
Algorithm~\ref{alg:loop} states the procedure; Figure~\ref{fig:schematic} shows one group.

\begin{figure*}[t]
\centering
\includegraphics[width=0.82\textwidth]{figures/method_grpo_group.png}
\caption{One GRPO group under look-ahead. The policy samples $G{=}8$ completions of a prompt; each
is rolled forward $K$ turns by the patient simulator $P$ and the current policy $\pi_n$; the oracle
scores the extended transcript; the scores are standardised within the group and every completion
carries a gradient. At $K{=}0$ the rollout is skipped, the oracle scores the candidate alone, and
standard GRPO is recovered: the horizon changes only what the oracle sees.}
\label{fig:schematic}
\end{figure*}

\begin{algorithm}[t]
\caption{\LAGRPO\ (one arm)}
\label{alg:loop}
\small
\begin{algorithmic}[1]
\REQUIRE base policy $\pi_1$; simulator $P$; oracle $O$; horizon $K$; group size $G$; MCL;
  iterations $N$
\FOR{$n = 1, \dots, N$}
  \STATE $\mathcal{C} \gets$ 96 conversations of $\pi_n$ with $P$ \COMMENT{eval set of $\pi_n$}
  \STATE $\mathcal{D} \gets$ patient-ending prefixes of $\mathcal{C}$ with $\geq$ MCL turns
  \FOR{each $c \in \mathcal{D}$ (two epochs)}
    \STATE $t_1, \dots, t_G \sim \pi(\cdot \mid c)$
    \STATE $r_g \gets O\big(c \oplus t_g \oplus \tau_K(c \oplus t_g)\big)$ for all $g$
    \STATE $A_g \gets (r_g - \operatorname{mean}_g r)/\operatorname{std}_g r$
    \STATE PPO-clipped step on $\{(c, t_g, A_g)\}_g$, KL to $\pi_n$
  \ENDFOR
  \STATE $\pi_{n+1} \gets \pi$
\ENDFOR
\STATE 96 conversations of $\pi_{N+1}$ with $P$ \COMMENT{eval set of $\pi_{N+1}$}
\end{algorithmic}
\end{algorithm}

\textbf{The iterative loop.} At iteration $n$, $\pi_n$ holds 96 conversations with $P$, one per
patient persona; each is sliced after every patient turn whose prefix has at least MCL utterances
(the \emph{minimum context length}, 12 here; see below) to form the training prompts. For each
prompt GRPO samples $G{=}8$ completions, scores them, standardises the
scores within the group, and takes a PPO-clipped step with a KL penalty toward $\pi_n$; after two
epochs the policy is saved as $\pi_{n+1}$ and the loop repeats with fresh conversations. The
conversations generated at the start of iteration $n$ double as the evaluation set of $\pi_n$; a
generate-only pass after the last update evaluates the final policy, so ten iterations yield
eleven evaluated model states per run.

\textbf{The look-ahead reward.} The only thing we change is the transcript the oracle scores.
Let $\tau_K(x)$ be the $K$-turn rollout of a transcript $x$: $P$ replies, then $\pi_n$, then $P$,
and so on for $K$ turns, stopping early if the patient closes the session; $\tau_0(x)$ is empty.
Then
\[
r_g \;=\; O\big(c \oplus t_g \oplus \tau_K(c \oplus t_g)\big).
\]
At $K{=}5$ the oracle scores the candidate together with the five turns that followed it under the
current policy, so a candidate is credited for where it leads under the therapist that will
actually follow it. Loss, advantage, KL penalty, prompts and number of gradient steps are identical
across $K$; the horizon changes only the evidence on which a candidate is ranked against its
siblings.

\textbf{Why the transfer is not trivial.} A rollout is one sample from the distribution of
continuations, so $r_g$ is a one-sample Monte Carlo estimate of the candidate's downstream value,
with sampling noise from $P$ and $\pi_n$ added to the oracle's own. PTO filters that noise by
construction, keeping only the best- and worst-scoring candidates, so an unlucky rollout can at
most cost a pair; GRPO standardises the scores within the group and trains on all $G$ candidates,
so the rollout variance enters every advantage. Whether the lever survives that exposure is the
question the controlled pair answers.

\textbf{Minimum context length.} Training prompts are prefixes, so the oracle scores a partial
conversation while the evaluation scores a whole session. In a pilot on an earlier configuration
of this task the oracle's ranking of very short prefixes agreed poorly with its ranking of the
completed sessions and recovered from roughly ten utterances, so both arms discard prefixes
shorter than MCL${=}12$. At that minimum, in the runs reported here, the training reward orders
candidates as the full-session score does in $86$--$89\%$ of pairs under the training oracle, and
agreement does not fall with prefix length (Appendix~\ref{app:mechanism}).

\textbf{Cost.} Each $K{=}5$ reward adds three patient-simulator calls and two policy generations
to the single oracle call both arms share, and roughly doubles per-step wall-clock (Limitations).
All comparisons are at matched iteration count, not matched cost.
